% ===========================================================================
%  离散数学与结构（I） 第二周（上）课堂笔记
%  代入的规定 · 一阶语义 · 量词与等号规则 · 亨金完备性 · Löwenheim–Skolem
%  Compile:  latexmk -xelatex main      (or: xelatex main, twice)
% ===========================================================================
\documentclass[cjk]{loom}
\usepackage{bussproofs}

\newcommand{\sem}[1]{[\![\,#1\,]\!]}        % 解释/真值记号
\newcommand{\N}{\mathbb{N}}
\newcommand{\Z}{\mathbb{Z}}
\newcommand{\Q}{\mathbb{Q}}
\newcommand{\R}{\mathbb{R}}
\newcommand{\M}{\mathcal{M}}
\newcommand{\Free}{\mathrm{Free}}
\newcommand{\PA}{\mathsf{PA}}

\runningthread{离散数学与结构 · 第二周（上）}

\begin{document}

\loomcover
  {一阶逻辑：代入、语义与完备性}
  {离散数学与结构（I）· 第二周（上）课堂笔记}
  {}
  {2026 年 9 月}

\begin{strand}
本讲内容为：代入的规定与变量捕获；一阶语言的语义（结构、解释、满足与模型）；
量词与等号的推演规则，并以 Peano 公理为例；完备性定理的亨金（Henkin）证明；
以及由完备性推出的 Löwenheim--Skolem 定理与斯科伦悖论。
\end{strand}

\clearpage
{\small\tableofcontents}

\clearpage
% ===========================================================================
\section{代入的规定}\label{sec:subst}

代入 $\varphi[t/x]$ 是用项 $t$ 替换公式 $\varphi$ 中 $x$ 的所有自由出现。
代入应当保持公式的语义：代入后的公式表达的应是“把 $x$ 所指的对象换成 $t$ 所指的对象
之后，原来的陈述”。唯一需要排除的情形是\keyword{变量捕获}：$t$ 中自由出现的变量
落入了原公式某个量词的作用域，从而被该量词约束，公式的含义随之改变。

\block{封闭项与句子}
不含自由变元的公式与项统称为\keyword{封闭的}（closed）：
\begin{itemize}[leftmargin=1.3em,itemsep=1pt]
  \item \keyword{封闭公式}（closed formula），即无自由变元的公式，称为\keyword{句子}
        （sentence）。句子的真值不依赖对变量的赋值，在一个结构中或真或假；
  \item \keyword{封闭项}（closed term），即无变元的项，在结构中解释为确定的论域元素。
\end{itemize}

\begin{definition}[可自由代入]
设 $\varphi$ 为公式，$x$ 为变元，$t$ 为项。称 $t$ 在 $\varphi$ 中\emph{可自由代入} $x$
（$t$ is free for $x$ in $\varphi$），若 $x$ 的每个被替入的自由出现位置都不处于
任何 $Qy$ 的作用域之内，其中 $y\in\operatorname{Var}(t)$ 且 $Q\in\{\forall,\exists\}$。
\end{definition}

该条件是对实际替入位置的局部条件：不要求 $t$ 的变元与 $\varphi$ 的全部约束变元
互不相同，只要求替入之后 $t$ 的变元不被捕获。若 $x$ 在 $\varphi$ 中没有自由出现，
代入不改变公式。

\begin{example}[变量捕获]\label{ex:capture}
设 $\psi=\exists x\,(x=y)$，考察 $\psi[(x+1)/y]$。
\end{example}

直接机械替换得到 $\exists x\,(x=x+1)$。$\psi$ 本来的含义是“存在与 $y$ 相等的对象”，
在论域非空的结构中恒真；替换结果却表示“存在等于自身后继的对象”——
项 $x+1$ 中自由出现的 $x$ 被 $\exists x$ 捕获。正确的做法是先对约束变元\keyword{改名}，
取公式中未出现的新变元 $u$：
\[
  \psi \;\equiv\; \exists u\,(u=y)
  \quad\Longrightarrow\quad
  \psi[(x+1)/y]=\exists u\,(u=x+1).
\]
此时 $x$ 仍然自由，语义保持为“存在与 $x+1$ 相等的对象”。

\begin{remark}
约束变元的名称不影响公式的含义：对新变元 $u$，$\exists x\,\varphi$ 与
$\exists u\,\varphi[u/x]$ 语义相同。因此先改名、后代入，可以把任意代入化为
合法代入。
\end{remark}

% ===========================================================================
\section{结构与满足}\label{sec:semantics}

语法规定符号串的合法构造，语义则规定符号的解释。一阶语言中的函数符号、
谓词符号与常元符号本身只是名称，其含义由结构指定。

\begin{definition}[结构]
一阶语言的一个\emph{结构} $\M=(\Omega;\ \ldots)$ 由以下部分组成：
\begin{itemize}[leftmargin=1.3em,itemsep=1pt]
  \item \keyword{论域} $\Omega$：一个\emph{非空}集合（universe）；
  \item 对每个 $k$ 元函数符号 $f_i^k$，指定函数 $f_i^k:\Omega^k\to\Omega$；
  \item 对每个 $k$ 元谓词符号 $P_i^k$，指定关系 $P_i^k\subseteq\Omega^k$；
  \item 对每个常元符号 $c_i$，指定元素 $c_i\in\Omega$。
\end{itemize}
\end{definition}

\begin{remark}[论域非空的理由]
若允许 $\Omega=\varnothing$，则 $\forall x\,(x=x)$ 在空论域上空虚地为真，
而 $\exists x\,(x=x)$ 为假。相应的推理链
\[
  \forall x\,(x=x)\ \longrightarrow\ x=x\ \longrightarrow\ \exists x\,(x=x)
\]
在空论域上失效：第一步是全称消去，第二步是存在引入，而后者要求论域中
确实存在可作见证的元素。即由全称命题推不出相应的存在命题。
允许空论域的逻辑体系称为\emph{自由逻辑}（free logic）；本课程一律假定论域非空。
\end{remark}

\block{解释的递归定义}
给定结构 $\M$，项的解释与公式的真值 $\sem{\cdot}_{\M}$ 递归定义如下：
\begin{align*}
  \sem{c_i}_{\M} &= c_i,\\
  \sem{f_i^k(t_1,\ldots,t_k)}_{\M} &= f_i^k\bigl(\sem{t_1}_{\M},\ldots,\sem{t_k}_{\M}\bigr),\\
  \sem{P_i^k(t_1,\ldots,t_k)}_{\M} &= 1
    \ \Longleftrightarrow\ \bigl(\sem{t_1}_{\M},\ldots,\sem{t_k}_{\M}\bigr)\in P_i^k,\\
  \sem{\forall x\,\psi}_{\M} &= \min_{v\in\Omega}\ \sem{\psi[v/x]}_{\M}.
\end{align*}
末行右端的 $v/x$ 需要说明：$v$ 是论域中的元素，并非语言中的符号，直接代入
不合语法。处理方式是\keyword{扩张语言}：为每个 $v\in\Omega$ 引入一个以它为解释的
常元符号，从而 $\psi[v/x]$ 仍是合法代入。等价的常见处理是使用变量赋值
$s[x\mapsto v]$。

\begin{remark}[等号的解释]
等号不作为普通的二元谓词符号处理，其解释固定为论域上的相等关系：
\[
  \sem{t_1=t_2}_{\M}=1 \ \Longleftrightarrow\ \sem{t_1}_{\M}=\sem{t_2}_{\M}.
\]
\end{remark}

\begin{definition}[满足、模型、逻辑后承]\label{def:satisfaction}
\leavevmode
\begin{itemize}[leftmargin=1.3em,itemsep=2pt]
  \item 对\emph{句子} $\varphi$：$\M\models\varphi$ 当且仅当 $\sem{\varphi}_{\M}=1$；
        此时称 $\M$ 为 $\varphi$ 的一个\keyword{模型}（model）。
  \item 对一般\emph{公式} $\varphi$（自由变元为 $x_1,\ldots,x_n$）：
        $\M\models\varphi$ 定义为 $\M\models\forall x_1\forall x_2\cdots\forall x_n\,\varphi$，
        即取\emph{全称闭包}。按此定义，$\M\models\varphi$ 同样不依赖变量赋值。
  \item 对公式集 $\Gamma$：$\M\models\Gamma$ 表示 $\M\models\varphi$ 对所有 $\varphi\in\Gamma$ 成立。
  \item $\models\varphi$（$\varphi$ \emph{有效}）表示 $\M\models\varphi$ 对所有结构 $\M$ 成立。
  \item $T\models\varphi$（\emph{逻辑后承}）表示：对所有结构 $\M$，$\M\models T$ 蕴含 $\M\models\varphi$。
\end{itemize}
\end{definition}

句子与一般公式分开定义的原因：句子的真值本就不依赖变量赋值，而含自由变元的
公式须先以全称量词约束全部自由变元，满足关系才有确定含义。

% ===========================================================================
\section{量词规则与等号规则}\label{sec:rules}

一阶逻辑的自然演绎在命题逻辑规则之外增加量词与等号两组规则。
两个量词规则的附加条件分别对应前两节的内容：$\forall I$ 的条件保证变元的
任意性，$\forall E$ 的条件即定义中的可自由代入条件。

\block{全称量词}
\begin{center}
\begin{minipage}{.42\linewidth}\centering
\AxiomC{$\varphi$}
\RightLabel{$\forall I$}
\UnaryInfC{$\forall x\,\varphi$}
\DisplayProof
\end{minipage}\hfill
\begin{minipage}{.52\linewidth}\centering
\AxiomC{$\forall x\,\varphi$}
\RightLabel{$\forall E$}
\UnaryInfC{$\varphi[t/x]$}
\DisplayProof
\end{minipage}
\end{center}
以前提集记法写出附加条件：
\begin{itemize}[leftmargin=1.3em,itemsep=1pt]
  \item $\forall I$：若 $\Gamma\vdash\varphi$，且 $x\notin\Free(\psi)$ 对所有 $\psi\in\Gamma$
        成立，则 $\Gamma\vdash\forall x\,\varphi$。若 $x$ 在某个未解除的假设中自由出现，
        它便只是满足该假设的特定对象，不能据以引入全称。
  \item $\forall E$：若 $\Gamma\vdash\forall x\,\varphi$，且 $t$ 在 $\varphi$ 中可自由代入 $x$，
        则 $\Gamma\vdash\varphi[t/x]$。此条件与§\ref{sec:subst} 的定义一致，
        用以排除变量捕获。
\end{itemize}

\block{等号规则 RE1--RE4}
\begin{center}
\begin{minipage}{.22\linewidth}\centering
\AxiomC{}
\RightLabel{$\mathrm{RE}_1$}
\UnaryInfC{$t=t$}
\DisplayProof
\end{minipage}\hfill
\begin{minipage}{.24\linewidth}\centering
\AxiomC{$t_1=t_2$}
\RightLabel{$\mathrm{RE}_2$}
\UnaryInfC{$t_2=t_1$}
\DisplayProof
\end{minipage}\hfill
\begin{minipage}{.42\linewidth}\centering
\AxiomC{$t_1=t_2$}
\AxiomC{$t_2=t_3$}
\RightLabel{$\mathrm{RE}_3$}
\BinaryInfC{$t_1=t_3$}
\DisplayProof
\end{minipage}
\end{center}
\begin{proposition}[$\mathrm{RE}_4$，代换规则]
相等的项可以在函数与谓词中逐项替换：
\[
  \frac{t_i=t_i'}{f(t_1,\ldots,t_i,\ldots,t_n)=f(t_1,\ldots,t_i',\ldots,t_n)},
  \qquad
  \frac{t_i=t_i'\qquad P(t_1,\ldots,t_i,\ldots,t_n)}{P(t_1,\ldots,t_i',\ldots,t_n)}.
\]
\end{proposition}
$\mathrm{RE}_1$ 至 $\mathrm{RE}_3$ 仅使 $=$ 成为等价关系；$\mathrm{RE}_4$ 使其与
所有函数、谓词相容，与语义中“等号固定解释为相等”的规定相一致。
这组规则也可以公理模式的形式给出，例如 $\mathrm{RE}_1$ 对应 $\forall x\,(x=x)$。

\begin{example}[$\exists x\,\forall y\,\varphi(x,y)\vdash\forall y\,\exists x\,\varphi(x,y)$]
将 $\exists$ 按 $\neg\forall\neg$ 展开，即证
$\neg\forall x\,\neg\forall y\,\varphi(x,y)\vdash\forall y\,\neg\forall x\,\neg\varphi(x,y)$：
\end{example}
\begin{center}
\AxiomC{$[\forall x\,\neg\varphi(x,y)]^1$}
\RightLabel{$\forall E$}
\UnaryInfC{$\neg\varphi(x,y)$}
\AxiomC{$[\forall y\,\varphi(x,y)]^2$}
\RightLabel{$\forall E$}
\UnaryInfC{$\varphi(x,y)$}
\RightLabel{$\neg I$（得 $\bot$）}
\BinaryInfC{$\bot$}
\RightLabel{$\to I_2$}
\UnaryInfC{$\neg\forall y\,\varphi(x,y)$}
\RightLabel{$\forall I$}
\UnaryInfC{$\forall x\,\neg\forall y\,\varphi(x,y)$}
\AxiomC{$\neg\forall x\,\neg\forall y\,\varphi(x,y)$}
\RightLabel{$\neg I$}
\BinaryInfC{$\bot$}
\RightLabel{$\to I_1$}
\UnaryInfC{$\neg\forall x\,\neg\varphi(x,y)$}
\RightLabel{$\forall I$}
\UnaryInfC{$\forall y\,\neg\forall x\,\neg\varphi(x,y)$}
\DisplayProof
\end{center}
两次 $\forall I$ 的附加条件均需检验：第一次施用时尚未解除的假设为前提与假设 $1$，
$x$ 在前提中无自由出现，在假设 $1$ 中被其自身的 $\forall x$ 约束；
最后一次施用时唯一未解除的假设为前提，$y$ 在其中无自由出现。两处均合法。

反向的 $\forall y\,\exists x\,\varphi\vdash\exists x\,\forall y\,\varphi$ 不可证。
取 $\varphi(x,y)$ 为 $x=y$，论域为任意至少含两个元素的集合：对每个 $y$ 都存在
与之相等的 $x$（取 $x=y$），故 $\forall y\,\exists x\,(x=y)$ 在该结构中为真；
但不存在固定的 $x$ 等于所有 $y$，故 $\exists x\,\forall y\,(x=y)$ 为假。
直观地说，$\exists x\,\forall y$ 要求同一个 $x$ 对所有的 $y$ 成立，
而 $\forall y\,\exists x$ 中的 $x$ 允许随 $y$ 变动。

\block{Peano 公理}
以下句子集记为 $\PA$，是后面讨论模型时的主要例子。语言含常元 $0$、
一元函数符号 $S$ 与二元函数符号 $+,\cdot$；论域的预期解释为 $\N$：
\begin{align*}
  &\forall x\ \neg\bigl(0=S(x)\bigr),
  &&\forall x\,\forall y\ \bigl(S(x)=S(y)\to x=y\bigr),\\
  &\forall x\ x+0=x,
  &&\forall x\,\forall y\ x+S(y)=S(x+y),\\
  &\forall x\ x\cdot 0=0,
  &&\forall x\,\forall y\ x\cdot S(y)=x\cdot y+x,
\end{align*}
另加\keyword{归纳公理模式}：对\emph{每个}公式 $\varphi$ 各取一条公理
\[
  \bigl(\varphi(0)\wedge\forall x\,(\varphi(x)\to\varphi(S(x)))\bigr)\longrightarrow\forall x\,\varphi(x).
\]
归纳并非一条公理，而是无穷多条公理组成的模式：一阶语言中不允许出现
$\forall\varphi$，即不能直接对“性质”作量化。这一限制是§\ref{sec:ls} 中
“一阶公理无法刻画 $\N$”的原因之一。

% ===========================================================================
\section{完备性定理：亨金构造}\label{sec:henkin}

\begin{strand}
完备性定理断言 $\Gamma\models\varphi\Rightarrow\Gamma\vdash\varphi$：
凡被所有模型满足的，都可以形式推导出来。对 $\Gamma\cup\{\neg\varphi\}$ 取逆否，
定理化为一个更对称的形式——
\[
  \Gamma\nvdash\bot \quad\Longleftrightarrow\quad \Gamma\text{ 有模型},
\]
即\emph{一致的句子集必有模型}。从右到左是可靠性，对推导树作归纳即可；
本节处理从左到右。
\end{strand}

\block{困难所在}
一致性是纯语法性质，它只说不存在以 $\Gamma$ 为前提、以 $\bot$ 为结论的有限推导；
而一个模型需要两样语义对象：一组充当论域的元素，以及每个原子句子在其上的真值。
直接面对一般的 $\Gamma$，两者都没有着落。具体有两处缺口：
\begin{enumerate}[label=(\alph*),leftmargin=1.8em,itemsep=1pt]
  \item \emph{见证缺口}。若 $\exists x\,\varphi(x)$ 可由 $\Gamma$ 推出，模型中就应当有
        具体的元素满足 $\varphi$。亨金模型打算用闭项（其等价类）充当论域元素，
        但若语言中根本没有为此准备的常元，存在句子的真值就没有着落。
  \item \emph{真值缺口}。$\Gamma$ 一般不能判定每个句子：对许多 $\varphi$，
        既有 $\Gamma\nvdash\varphi$ 又有 $\Gamma\nvdash\neg\varphi$。
        以“是否属于 $\Gamma$”来裁定真值，对这样的 $\varphi$ 没有定义。
\end{enumerate}
亨金构造依次补上这两个缺口：先向语言中加入见证常元，再把理论扩充为极大一致集。

\begin{definition}[理论、亨金系统]
给定公理集 $A$，其\emph{理论}为 $T=\{\text{句子 }\varphi\mid A\vdash\varphi\}$。
称 $T$ 为\keyword{亨金的一阶系统}（Henkin theory），若对每个形为
$\exists x\,\varphi(x)$ 的句子，都存在常元符号 $c$（称为\emph{亨金见证}，
Henkin witness），使
\[
  \bigl(\exists x\,\varphi(x)\bigr)\longrightarrow\varphi(c)\ \in\ T.
\]
换言之，理论一旦断言某个性质有实例，语言中就备好了指称该实例的常元。
\end{definition}

\begin{theorem}[模型存在定理]\label{thm:completeness}
在上述含量词规则与等号规则的自然演绎系统中，一致的句子集必有模型。
由此即得完备性：$\Gamma\models\varphi$ 蕴含 $\Gamma\vdash\varphi$。
\end{theorem}

\begin{proof}
分三步。

\par\smallskip\noindent\textbf{第一步：亨金扩充（Henkin extension）}\quad
原理论一般不是亨金的，因此造递增链 $T=T_0\subseteq T_1\subseteq T_2\subseteq\cdots$：
$T_{i+1}$ 在 $T_i$ 的语言中对每个句子 $\exists x\,\varphi(x)$ 引入一个\emph{新}常元
$c_\varphi$，并加入公理
\[
  \exists x\,\varphi(x)\longrightarrow\varphi(c_\varphi).
\]
必须逐层迭代而不能一次完成：新常元进入语言后会产生新的存在句子，它们又需要
各自的见证。取 $T'=\bigcup_{i\ge0}T_i$，则 $T'$ 是亨金的。
一致性在扩充下保持：若某个 $T_i\vdash\bot$，推导只涉及有限多条新见证公理；
从新到旧逐条把这些公理还原为相应的存在句子（$\exists$ 消去与引入），
便在 $T$ 中复现了矛盾，与 $T$ 一致矛盾。

\par\smallskip\noindent\textbf{第二步：极大一致扩充}\quad
由 Lindenbaum 引理，$T'$ 包含于某个极大一致集 $T_{\max}$。该引理第一周已在
命题逻辑情形证明；一阶情形的论证相同：枚举全部句子，逐个吸收且保持一致，
推导的有限性保证并集仍一致。极大一致性给出两条关键性质：每个句子 $\varphi$
与 $\neg\varphi$ 恰有一个属于 $T_{\max}$；且 $T_{\max}$ 对推导封闭。
真值缺口由此补上：每个句子的“真值”现在都有定义——它是否在 $T_{\max}$ 中。

\par\smallskip\noindent\textbf{第三步：构造模型}\quad
论域取亨金常元的\keyword{等价类}：
\[
  \Omega=\bigl\{[c_\varphi]\ \big|\ \exists x\,\varphi(x)\in T_{\max}\bigr\},
  \qquad
  c_\varphi\sim c_\psi \ \Longleftrightarrow\ (c_\varphi=c_\psi)\in T_{\max}.
\]
由 RE1--RE3，$\sim$ 是等价关系。函数与谓词按语法方式解释：
$f^{\M}([c_1],\ldots,[c_k])=[f(c_1,\ldots,c_k)]$，以及
$([c_1],\ldots,[c_k])\in P^{\M}$ 当且仅当 $P(c_1,\ldots,c_k)\in T_{\max}$；
解释的良定义性（与代表元选取无关）由代换规则 $\mathrm{RE}_4$ 保证。
至此每个原子句子的真值都已确定。对公式作结构归纳，证明\emph{真值引理}
\[
  \M\models\varphi \quad\Longleftrightarrow\quad \varphi\in T_{\max}.
\]
原子情形由定义；联结词情形由 $T_{\max}$ 的封闭性与完备性（$\varphi$ 与
$\neg\varphi$ 恰居其一）。关键是存在量词情形：若 $\exists x\,\varphi(x)\in T_{\max}$，
见证公理给出 $\varphi(c_\varphi)\in T_{\max}$，由归纳假设得
$\M\models\varphi(c_\varphi)$，故 $\M\models\exists x\,\varphi(x)$；
反之若 $\exists x\,\varphi(x)\notin T_{\max}$，则 $\neg\exists x\,\varphi(x)\in T_{\max}$，
由归纳假设与封闭性，不存在满足 $\varphi$ 的论域元素，故
$\M\not\models\exists x\,\varphi(x)$。这正是第一步加入见证常元的原因：
没有它们，$\M$ 中可能没有元素来实现理论断言的存在性。

于是 $\M\models T_{\max}\supseteq\Gamma$，即一致的 $\Gamma$ 有模型。
\end{proof}

\begin{corollary}[紧致性定理]
若句子集 $\Gamma$ 的每个有限子集都有模型，则 $\Gamma$ 有模型。
\end{corollary}
\begin{proof}
形式推导只含有限条前提，故“每个有限子集都有模型”经可靠性给出
“每个有限子集都一致”，进而 $\Gamma$ 一致；再由定理~\ref{thm:completeness} 得证。
\end{proof}

\begin{remark}\label{rem:countable}
亨金模型的论域由语言中的常元之等价类构成。语言至多可数时，常元至多可数，
故模型至多可数。下一节将用到这一观察。
\end{remark}

% ===========================================================================
\section{Löwenheim--Skolem 定理与斯科伦悖论}\label{sec:ls}

完备性定理与紧致性定理直接导致两个表面上与直觉相悖的结论：
\emph{自然数理论有不可数模型，实数理论有可数模型}。二者合起来表明，
一阶公理无法控制其模型的基数。

\begin{theorem}[上行：Peano 公理有不可数模型]\label{thm:upward}
存在 $\M\models\PA$ 使 $|\M|>|\N|$。
\end{theorem}
\begin{proof}
任取不可数集合 $S$。扩充语言：对每个 $s\in S$ 引入常元 $c_s$；
扩充公理：
\[
  c_s\neq c_{s'}\qquad(s\neq s',\ s,s'\in S).
\]
记所得句子集为 $\PA^{+}$。$\PA^{+}$ 的任意有限子集只涉及有限多个 $c_s$，
在 $\N$ 中将这有限个常元解释为互不相同的自然数，即得该有限子集的模型。
由紧致性定理，$\PA^{+}$ 有模型 $\M$。（等价地：矛盾推导只涉及有限条公理，
有限部分一致故 $\PA^{+}$ 一致，再应用定理~\ref{thm:completeness}。）
诸 $c_s^{\M}$ 两两不同，故 $|\M|\ge|S|>\aleph_0$。
\end{proof}

此类 $\M$ 中除标准元素 $0,S(0),S(S(0)),\ldots$ 外还含非标准元素——
例如任何 $c_s^{\M}$ 都不能由闭项 $S(\cdots S(0)\cdots)$ 表出。

\begin{theorem}[下行：可数语言的一致理论有可数模型]\label{thm:downward}
设语言至多可数，$T$ 一致，则 $T$ 有至多可数的模型。
\end{theorem}
\begin{proof}
由注~\ref{rem:countable}，定理~\ref{thm:completeness} 所构造的亨金模型
在语言至多可数时至多可数。应用于实数：实数的一阶理论 $T$ 以
$(\R;0,1,+,\cdot,<)$ 为模型，故一致，从而存在至多可数的 $\M\models T$。
换言之，\emph{实数有可数解释}。
\end{proof}

\begin{remark}[斯科伦悖论]
将下行定理用于实数理论（集合论情形类似），表面上出现矛盾：$\M$ 可数，
而 Cantor 对角线论证可在系统内形式化，因此“$\R$ 不可数”这句一阶句子
在 $\M$ 中为真。矛盾的错觉源于混淆两个层面：
\begin{itemize}[leftmargin=1.3em,itemsep=1pt]
  \item \keyword{系统内}：“不存在从 $\N$ 到 $\R$ 的满射”是 $\M$ 所满足的一阶句子，
        其含义仅为“$\M$ 中不存在这样的函数对象”；
  \item \keyword{系统外}：从外部考察，$\M$ 的论域确实可数，$\N$ 与 $\M$ 所解释的
        “实数”之间确实存在双射——但该双射是系统外的对象，语言中没有符号
        可以指称它，系统内也就无法陈述“此双射存在”。
\end{itemize}
因此，可数与不可数并非一阶语言能够刻画的绝对性质，而只是相对于模型而言的性质。
一阶公理既无法排除不可数的“自然数”，也无法排除可数的“实数”。
\end{remark}

\begin{remark}
若要唯一刻画 $\N$（在同构意义下），需要超出一阶逻辑的手段，例如二阶语义
（对论域的所有子集作量化）或 $\omega$-规则。另一方面，非标准模型并非只是障碍：
非标准分析正是以 $\R$ 的含无穷小元素的扩充模型为工具建立的。
\end{remark}


\end{document}
